%\documentclass[handout]{beamer}
\documentclass{BeamerTemplate}


\newcounter{sauvegardeenumi}
\newcommand{\asuivre}{\setcounter{sauvegardeenumi}{\theenumi}}
\newcommand{\suite}{\setcounter{enumi}{\thesauvegardeenumi}}


%\usepackage[latin1]{inputenc}  %modifies preambule

\title{Inférence}

\date{}


\begin{document}

\begin{frame}

\titlepage

\end{frame}

\begin{frame}{Tu te mets combien?}
 \begin{table}[h!]
    \centering
    \renewcommand{\arraystretch}{1.4}
    \begin{tabular}{c|c|c|c|c}
        \textbf{$\mathbb{E}$ et $V$} & 
        \textbf{Lois + TCL} & 
        \textbf{Estimation} & 
        \textbf{IC} & 
        \textbf{Tests d'hypothèses} \\ \hline
        \href{Niveau 1}{\ref{ev1}} & \href{Niveau 1}{\ref{ech1}} & \href{Niveau 1}{\ref{est1}} & \href{Niveau 1}{\ref{ic1}} & \href{Niveau 1}{\ref{tests1}} \\
        \href{Niveau 2}{\ref{ev2}} & \href{Niveau 2}{\ref{ech2}} & \href{Niveau 2}{\ref{est2}} & \href{Niveau 2}{\ref{ic2}} & \href{Niveau 2}{\ref{tests2}} \\
        \href{Niveau 3}{\ref{ev3}} & \href{Niveau 3}{\ref{ech3}} & \href{Niveau 3}{\ref{est3}} & \href{Niveau 3}{\ref{ic3}} & \href{Niveau 3}{\ref{tests3}} \\
        \href{Niveau 4}{\ref{ev4}} & \href{Niveau 4}{\ref{ech4}} & \href{Niveau 4}{\ref{est4}} & \href{Niveau 4}{\ref{ic4}} & \href{Niveau 4}{\ref{tests4}} \\
        \href{Niveau 5}{\ref{ev5}} & \href{Niveau 5}{\ref{ech5}} & \href{Niveau 5}{\ref{est5}} & \href{Niveau 5}{\ref{ic5}} & \href{Niveau 5}{\ref{tests5}} \\
    \end{tabular}
    \label{tab:niveaux}
\end{table}
\end{frame}

\begin{frame}{Espérance et variance 1}\label{ev1}
\begin{center}
        Soit la variable aléatoire $X$, dont les deux résultats possibles sont $X=1$ avec probabilité $0,4$ et $X=2$ avec probabilité $0,6$. 

    \medskip
    Quelle est l'espérance de $X$?
\end{center}

\end{frame}


\begin{frame}{Espérance et variance 2}\label{ev2}
\begin{center}
        Vrai ou faux? Si faux, corrigez.
\medskip

L'espérance d'une variable aléatoire est invariante par translation, mais sa variance n'est pas invariante par translation.
\end{center}

\end{frame}


\begin{frame}{Espérance et variance 3}\label{ev3}
\begin{center}
      Soit $X$ et $Y$, deux variables aléatoires indépendantes, d'espérances respectives $\mu_x$ et $\mu_Y$, et de variances respectives $\sigma^2_X$ et $\sigma^2_Y$.
      \medskip

      Donnez les expressions de l'espérance et de la variance de $4X-2Y$.
\end{center}

\end{frame}

\begin{frame}{Espérance et variance 4}\label{ev4}
\begin{center}
      Parmi les expressions suivantes, laquelle est une expression de la variance pour une variable aléatoire $X$ continue?

      \begin{enumerate}
          \item $V(X)=\int_{-\infty}^{\infty}(x-\mu)^2dx$
          \item $V(X)=E[(X-\mu)^2]$
          \item $V(X)=\sum_{\mathbb{R}_X}(x-\mu)^2p(x)$
          \item $V(X)=E(X)-E(X)^2$
      \end{enumerate}
\end{center}

\end{frame}

\begin{frame}{Espérance et variance 5}\label{ev5}
\begin{center}
      Soit $X$, une variable aléatoire uniforme dans l'intervalle $[12, 15]$, dont la fonction de répartion est donnée par $F(x)=\frac{x-12}{3}$ si $x \in [12,15]$ et $F(x)=0$, sinon. 

      \medskip

      Vous avez 4 minutes pour calculer l'espérance et la variance de $X$.
\end{center}

\end{frame}

\begin{frame}{Lois échantillonnales + TCL 1}\label{ech1}
\begin{center}

Vrai ou faux?
\medskip

      La loi de Student de degrés de liberté $\nu$ a une d'espérance égale à $\nu$. Si faux, corrigez.
\end{center}

\end{frame}

\begin{frame}{Lois échantillonnales + TCL 2}\label{ech2}
\begin{center}
Soit $X_1,\hdots, X_{n}$, des réalisation d'une variable aléatoire de loi quelconque d'espérance $\mu$ et de variance $\sigma^2$.

\medskip
Énoncez le théorème limite centrale pour la somme des $X_i$ ainsi que la condition d'application.

\end{center}

\end{frame}

\begin{frame}{Lois échantillonnales + TCL 3}\label{ech3}
\begin{center}

Vrai ou faux? Si faux, corrigez.
\medskip

Soit $X_1,\hdots, X_{35}$, des réalisation d'une variable aléatoire de loi quelconque d'espérance $\mu$ et de variance $\sigma^2$. Alors la variable aléatoire $\displaystyle \frac{\overline{X}-\mu}{S/\sqrt{35}}$ suit approximativement une loi de Student à 34 degrés de liberté.

\end{center}

\end{frame}

\begin{frame}{Lois échantillonnales + TCL 4}\label{ech4}
\begin{center}

Soit l'échantillon aléatoire $X_1,\hdots,X_{25}$ qui sont des réalisations d'une variable aléatoire $X$ qui suit une loi normale d'espérance $16$ et de variance $6400$.

\medskip

Quelle est la loi de la variable aléatoire $\displaystyle \frac{\overline{X}}{16}-1$?


\end{center}

\end{frame}


\begin{frame}{Lois échantillonnales + TCL 5}\label{ech5}
\begin{center}

Soit l'échantillon aléatoire $X_1,\hdots,X_{25}$ qui sont des réalisations d'une variable aléatoire $X$ qui suit une loi normale d'espérance $16$ et de variance $6400$.

\medskip

Quelle est la loi de la variable aléatoire $\left(\displaystyle \frac{\overline{X}}{16}-1\right)^2$?

\end{center}

\end{frame}

\begin{frame}{Estimation de paramètres 1}\label{est1}
\begin{center}

Quel est l'estimateur <<classique>> de l'espérance $\mu$?

\end{center}

\end{frame}


\begin{frame}{Estimation de paramètres 2}\label{est2}
\begin{center}

Pourquoi utilisons-nous $\displaystyle S^2=\frac{1}{n-1}\sum_{i=1}^n(X_i-\overline{X})^2$ plutôt que $\displaystyle S^2=\frac{1}{n}\sum_{i=1}^n(X_i-\overline{X})^2$ pour estimer $\sigma^2$?

\end{center}

\end{frame}

\begin{frame}{Estimation de paramètres 3}\label{est3}
\begin{center}
Soit $\overline{X}_1, \hdots , \overline{X}_m$, des réalisations de la moyenne empirique $\overline{X}$.

\medskip
Corrigez cette affirmation:

<<Comme $\overline{X}$ est estimateur sans biais pour $\mu$, nous savons que $\displaystyle \frac{1}{m}\sum_{i=1}^m \overline{X}_m=\mu$.>>

\end{center}

\end{frame}

\begin{frame}{Estimation de paramètres 4}\label{est4}
\begin{center}

L'erreur quadratique moyenne d'un estimateur sans biais est égale à....

\end{center}

\end{frame}

\begin{frame}{Estimation de paramètres 5}\label{est5}
\begin{center}

On estime le paramètre $\theta$ à l'aide de deux estimateurs :
\[
\begin{aligned}
\text{Estimateur 1:}\quad 
& \operatorname{Biais}_1(\theta)=\theta, 
& \operatorname{Var}_1(\theta)=12,\\[2mm]
\text{Estimateur 2:}\quad 
& \operatorname{Biais}_2(\theta)=\sqrt{\theta}, 
& \operatorname{Var}_2(\theta)=12.
\end{aligned}
\]

Pour \textbf{quelles valeurs de $\theta$} l’estimateur 1 est-il \textbf{préférable} selon le critère de l’erreur quadratique moyenne (EQM) ?

\end{center}
\end{frame}



\begin{frame}{Intervalle de confiance 1}\label{ic1}
\begin{center}

Quel énoncé parmi les suivant est vrai?

\medskip
\begin{enumerate}
    \item L'intervalle de confiance pour la variance est symétrique autour de $S^2$.
    \item L'intervalle de confiance pour l'espérance $\mu$ est symétrique autour de $\overline{X}$.
    \item Plus un intervalle de confiance est large, plus il est précis.
    \item La quantité $1-\alpha$ se nomme le seuil de l'intervalle de confiance.
\end{enumerate}

\end{center}

\end{frame}

\begin{frame}{Intervalle de confiance 2}\label{ic2}
\begin{center}

Comment se nomme la quantité $z_{\alpha/2}\frac{\sigma}{n}$ dans la construction d'un intervalle de confiance de niveau $1-\alpha$ pour l'espérance d'une loi normale avec variance connue?

\end{center}

\end{frame}

\begin{frame}{Intervalle de confiance 3}\label{ic3}
\begin{center}

Nous savons que $E(\widehat{p})=p$ et $\displaystyle V(\widehat{p})=\frac{p(1-p)}{n}$. Quel est le problème dans la construction de l'intervalle de niveau approximatif $1-\alpha$ pour $p$ suivant?

\begin{align*}
    \widehat{p}\pm z_{\alpha/2}\sqrt{\frac{p(1-p)}{n}}
\end{align*}

\end{center}

\end{frame}

\begin{frame}{Intervalle de confiance 4}\label{ic4}
\begin{center}

Soit l'intervalle de confiance $[L;U]$ de niveau $1-\alpha$ pour $\mu$, dans le cas de la loi normale, avec variance inconnue et construit avec un échantillon de taille $n$.
\medskip

Dites si chacun des énoncés suivants est vrai ou faux. Expliquez quand ils sont faux.

\begin{enumerate}
    \item La probabilité que $\overline{X}$ se trouve dans l'intervalle de confiance est $1-\alpha$.
    \item $L$ et $U$ sont des variables aléatoires car ils sont impossibles à calculer.
    \item Toutes choses étant égales par ailleurs, un intervalle de confiance construit avec la variance connue (et $\sigma^2=s^2$) est plus court.
    \item Si $[L;U]=[12,34;17,10]$, il nous manque d'information pour calculer la marge d'erreur.
    \item Un intervalle de confiance construit avec $m>>n$ aura un niveau de confiance plus élevé.
\end{enumerate}

\end{center}

\end{frame}

\begin{frame}{Intervalle de confiance 5}\label{ic5}
\begin{center}

Énoncez clairement TOUS les cas (et les conditions d'application) auxquels ces intervalles s'appliquent:

\begin{enumerate}
    \item $\displaystyle \overline{X}\pm t_{n-1;\alpha/2}\frac{S}{\sqrt{n}}$
    \item $\displaystyle \widehat{p}\pm z_{\alpha/2}\sqrt{\frac{\widehat{p}(1-\widehat{p})}{n}}$
    \item $\displaystyle \overline{X}\pm z_{\alpha/2}\frac{\sigma}{\sqrt{n}}$
    \item $\displaystyle \left[\frac{(n-1)S^2}{\chi^2_{n-1;\alpha/2}};\frac{(n-1)S^2}{\chi^2_{n-1;1-\alpha/2}}\right]$
    \item $\displaystyle \overline{X}\pm z_{\alpha/2}\frac{S}{\sqrt{n}}$
    
\end{enumerate}

\end{center}

\end{frame}

\begin{frame}{Tests d'hypothèse 1}\label{tests1}
\begin{center}

Quels sont les trois types d'hypothèses alternatives?

\end{center}

\end{frame}

\begin{frame}{Tests d'hypothèse 2}\label{tests2}
\begin{center}

Décrivez dans vos mots l'erreur de type I.

\end{center}

\end{frame}

\begin{frame}{Tests d'hypothèse 3}\label{tests3}
\begin{center}

Décrivez dans vos mots la puissance d'un test.

\end{center}

\end{frame}


\begin{frame}{Tests d'hypothèse 4}\label{tests4}
\begin{center}

Quelle est la condition nécessaire par rapport à l'échantillon dans le test sur la variance qui permet d'utiliser la statistique $\displaystyle \frac{(n-1)S^2}{\sigma^2}$ qui suit une $\chi_{n-1}^2$, sous $H_0$?

\end{center}

\end{frame}


\begin{frame}{Tests d'hypothèse 5}\label{tests5}
\begin{center}

Vous faites un test sur 20 observations  pour confronter les hypothèses $H_0: \mu=5$ et $H_1: \mu<5$. Vous obtenez un seuil observé (valeur P) de 0,09.

\medskip
Vrai ou faux? Vous avez 30 secondes par question.
\medskip

\begin{enumerate} \itemsep 3mm
\item La moyenne échantillonnale est inférieure à 5.
\item Vous rejetez $H_0$ au seuil de 5\%.
\item Vous rejetez $H_0$ au seuil de 10\%.
\item Vous rejetez $H_0: \mu=4$ contre $H_1: \mu<4$ au seuil de 5\%.
\item Vous refaites l'expérience avec 40 observations. Vous obtenez la même moyenne échantillonnale. Votre seuil observé est alors plus élevé que 0,09 pour le test $H_0: \mu=5$ contre $H_1: \mu<5$.
\end{enumerate}


\end{center}

\end{frame}



\end{document}


